% Dédicaces
\chapter*{Dédicaces}
\addcontentsline{toc}{chapter}{Dédicaces}

\noindent
Je dédie ce travail à ma famille pour son soutien, sa confiance et ses encouragements tout au long de mon parcours académique. Je le dédie également à mes enseignants et encadrants, dont les conseils et l'accompagnement ont contribué à développer mon esprit d'analyse, ma rigueur et mon autonomie. Enfin, je dédie ce rapport à toutes les personnes qui ont participé, directement ou indirectement, à la réussite de ce stage.

\newpage

% Remerciements
\chapter*{Remerciements}
\addcontentsline{toc}{chapter}{Remerciements}

\noindent
Je tiens à exprimer ma profonde gratitude à toutes les personnes qui ont contribué à la réussite de ce stage et à la réalisation de ce rapport. Je remercie particulièrement mon encadrant académique et l'équipe pédagogique de la filière \filiere{} pour leur suivi, leurs orientations et leur disponibilité.

\noindent
Je remercie également l'équipe de \entreprise{} pour son accueil, son accompagnement et les échanges constructifs durant la période du stage. Les informations administratives et les noms des encadrants en entreprise seront complétés après validation officielle par l'établissement d'accueil.

\newpage

% Résumé
\chapter*{Résumé}
\addcontentsline{toc}{chapter}{Résumé}

\noindent
Ce rapport présente le travail réalisé dans le cadre de mon stage autour du projet \projectName{}, une plateforme web de recrutement destinée au secteur CHR au Maroc, c'est-à-dire les cafés, hôtels et restaurants. L'objectif du projet est de centraliser le processus de recrutement, souvent informel dans ce secteur, en proposant un espace structuré pour les candidats, les recruteurs et l'administrateur.

\noindent
La plateforme permet aux candidats de créer un profil professionnel, d'ajouter leur CV une seule fois, de consulter les offres, de postuler et de passer un quiz métier lorsque l'offre l'exige. Les recruteurs disposent d'un tableau de bord pour créer des offres, ajouter éventuellement un quiz, consulter les candidatures, accéder aux CV et prendre une décision d'acceptation ou de refus. L'administrateur peut suivre les statistiques globales, gérer les utilisateurs et modérer les offres.

\noindent
La solution a été développée avec une architecture web moderne : \textbf{React}, \textbf{Vite}, \textbf{Tailwind CSS} et \textbf{Framer Motion} pour le frontend, \textbf{Laravel 11}, \textbf{Sanctum} et \textbf{MySQL} pour le backend et la persistance des données. Le travail réalisé couvre l'analyse des besoins, la conception, l'implémentation des modules principaux, les tests fonctionnels et la préparation d'un scénario de démonstration.

\newpage

\tableofcontents

\chapter*{Liste des abréviations}
\addcontentsline{toc}{chapter}{Liste des abréviations}

\begin{table}[h!]
\centering
\renewcommand{\arraystretch}{1.4}
\setlength{\tabcolsep}{18pt}
\begin{tabular}{|l|p{10cm}|}
\hline
\textbf{Abréviation} & \textbf{Signification} \\
\hline
API & Application Programming Interface \\
BD & Base de données \\
CHR & Café, Hôtel, Restaurant \\
CRUD & Create, Read, Update, Delete \\
CV & Curriculum Vitae \\
HTTP & HyperText Transfer Protocol \\
JSON & JavaScript Object Notation \\
QCM & Questionnaire à choix multiples \\
REST & Representational State Transfer \\
SaaS & Software as a Service \\
UML & Unified Modeling Language \\
UX & User Experience \\
Vite & Outil de build frontend pour applications web modernes \\
Laravel Sanctum & Système d'authentification par tokens pour API Laravel \\
\hline
\end{tabular}
\end{table}

\newpage

\listoffigures
\addcontentsline{toc}{chapter}{Liste des figures}

\newpage

\listoftables
\addcontentsline{toc}{chapter}{Liste des tableaux}
